% Owner: Role A. Translated and adapted from the verified Vietnamese draft
% docs/report_draft_g_h.md (section "h. Conclusion"), AFTER the 2026-08-30
% fact-check pass that fixed a fabricated rule citation (the Dropout rule
% quoted here was cross-checked against docs/report_draft_d_e.md's actual
% extracted rule, not invented -- see progress/A.md log). If Role B's
% e_analysis.tex rule wording changes, update the quoted rule below to
% match.

\section{Conclusion}
\label{sec:h}

\subsection{What the group learned}

The project directly demonstrates the two theoretical points raised in
the Introduction: a decision tree is both interpretable and prone to
overfitting if its complexity is not controlled. After nominal categories
are encoded, it also needs no feature scaling. The unconstrained M0 model
reaches perfect training accuracy (1.000) while test accuracy is only
0.6689, a generalization gap of 0.3311 measured directly from this
project's own results rather than assumed from theory
\citep{breiman1984classification, mitchell1997machine}.

The three improvement directions highlight distinct lessons:

\begin{enumerate}
  \item \textbf{For the severely overfit baseline in this project,
    controlling complexity is highly effective.} Pruning improves aggregate
    held-out performance and structural simplicity, although its effect is
    not uniform across classes because Enrolled recall decreases slightly.
  \item \textbf{Class-imbalance techniques do not automatically help.}
    Effectiveness depends on the specific mechanism (M2a reweights the
    Gini computation without adding data, while M2b adds synthetic
    samples in the minority class's feature region), and must be judged
    by per-class recall, not overall accuracy alone: a model can be
    ``worse'' on accuracy yet ``better'' on the real application goal.
  \item \textbf{The most accurate model is not always the most useful
    one.} M3 shows that this dataset's strongest predictive features
    (semester outcomes) only become available well after enrollment,
    the same window during which an early-warning model would need to
    act, so M3's lower accuracy is a reasonable price for being usable
    that early.
\end{enumerate}

\subsection{Key findings from the experiments}

\begin{itemize}
  \item M0's overfitting: a 0.3311 generalization gap (train 1.000 vs.\
    test 0.6689), depth 27, 634 leaves. The tree ``memorizes'' rather
    than generalizes.
  \item Pruning shrinks the leaf count by \textbf{$\sim$37$\times$}
    (634$\to$17) while raising test accuracy by 8.70 percentage points,
    consistent with much of the removed structure capturing sample-specific
    variation rather than stable predictive signal.
  \item \texttt{class\_weight='balanced'} (M2a) does not improve
    minority-class recall on this data. SMOTE-on-train (M2b) raises
    Enrolled recall by 8.18 points and Dropout recall by 2.82 points
    without a corresponding drop in Enrolled precision
    (Section~\ref{sec:g}), so the gain is not pure over-prediction of
    Enrolled -- though SMOTE on one-hot-encoded categorical features
    still carries known limitations.
  \item Dropping the 12 semester-outcome columns (M3) costs 12.77
    accuracy points, a direct measurement of the ``price'' of wanting
    an early warning instead of a later but more accurate one.
  \item All 5 models share one fixed, stratified train/test split
    (\texttt{random\_state=42}). This makes the comparisons directly
    comparable and prevents different partitions from confounding the
    reported differences, but it does not eliminate sampling uncertainty.
\end{itemize}

\subsection{Effectiveness of decision trees for this problem}

A decision tree suits this dataset for the reasons argued in the Dataset
Description section: features have clear semantics, so every decision
rule is interpretable. For example, the Dropout rule extracted from the
M0 tree in Section~\ref{sec:e} is a chain of 9 conditions that together
identify a leaf of $n=189$ training samples at 100\% purity. The single
most readable condition in that chain is passing at most 1.5
second-semester course units (\texttt{2nd sem approved} $\le$ 1.5), but
that threshold alone is not sufficient by itself; it is one term in the
full conjunction. The parallel Graduate leaf ($n=176$, 100\% purity) is
a 16-condition chain, the two most readable conditions being passing
more than 5.5 second-semester course units \emph{and}
\texttt{Tuition fees up to date} $> 0.5$, consistent with the EDA
observation in Figure~\ref{fig:c2} that paying tuition on time
correlates strongly with graduating. Long condition chains like these
are themselves a sign of the overfitting discussed above: a rule that
needs 9--16 ANDed conditions to stay 100\% pure is fitting narrow
pockets of the training set rather than a simple general pattern, one of
the concrete costs of leaving M0 unpruned. After preprocessing, the tree
can split numerical and one-hot indicator columns without feature
scaling. But the project also exposes the inherent limits of a single,
unconstrained tree: it overfits easily, and no single model
configuration is simultaneously the most accurate, the most balanced
across classes' recall, and the earliest-usable; these three goals
require three different configurations of the same algorithm. (This is
a class-wise recall comparison, not a formal fairness audit across
sensitive attributes such as gender or nationality, which this project
does not perform.) For an application with multiple distinct goals, like
dropout prediction, a more realistic conclusion than ``one best model''
is that \textbf{a decision tree is a flexible enough tool to be
optimized for each goal separately} (accuracy, class-wise recall
balance, or timeliness), provided the user chooses the right
configuration and understands the trade-off being accepted. This is the
kind of explanation the assignment brief asks for: why each choice helps
or does not, not only which one scores highest.
